\cotefile{10}
\scanpage{2}

\uline{Polarisations sur les} \uline{$\otimes$-catégories \unsure{galoisiennes|gérinales|graduées}}

% archive sheet number: 1
\scanpage{3}

% The hand and its underline run upside down in the supplied scan.
\uline{Photocopies}
\scanpage{4}

\deleted{$\mathcal M$ catégorie abélienne $K$-linéaire, $K\subset\mathbb R$}

\deleted{\unsure{a)}} $\mathcal M$ une $\otimes$-catégorie sur $K\subset\mathbb R$, avec \deleted{\unsure{graduation et}} \inserted{\deleted{\unsure{soit}} un} objet inversible \deleted{$K(1)$ de degré $2$.} $L$

1) Soit $M\in\operatorname{Ob}\mathcal M^{\unsure{i|*}}$, $\varphi:M\otimes M\to L\deleted{\unsure{K(1)}}$. Supposons $\varphi$ symétrique ou alternée, $\varphi$ non dégénérée. Alors $\varphi$ définit sur $\operatorname{End}(M)$ une \uline{anti-involution} $u\mapsto u^*$ $[\,\varphi(ux,y)=\varphi(x,u^*y)\,;$ posant $\varphi(x,y)=\langle x,\widetilde\varphi y\rangle$, \inserted{$\widetilde\varphi:M\to\check M\otimes L$,} on aura donc $u^*=\widetilde\varphi^{-1}\circ\check u\otimes1\circ\widetilde\varphi\,]$. \deleted{\unsure{Tout}} On a une bijection

\[
\operatorname{Hom}(M\otimes M,L\deleted{\unsure{K(1)}})\simeq\operatorname{End}(M)
\]

[$\psi$ et $u$ se correspondent si $\psi(x,y)=\varphi(x,uy)$, i.e $\widetilde\psi=\widetilde\varphi u$], et cette dernière \deleted{\unsure{oblige|bijection}} transforme $\psi\mapsto\varepsilon{}^{\unsure{s|t}}\psi$ en $u\mapsto\varepsilon u^*$, (où $\varepsilon=\pm1$ suivant que $\varphi$ sym. ou alternée). Donc les formes sym (resp. alternées) de $\operatorname{Hom}(M\otimes M,L)$ correspondent aux $u$ tels que $u^*=u$.

2) On dit que $\varphi$ \deleted{\unsure{est p|est une|est définie}} satisfait la \uline{condition de Weil} si \inserted{$\varphi$ (sym. ou \unsure{antisym.|altern.} non dég) est telle que} $\operatorname{Tr}u^*u>0$ pour $u\in\operatorname{End}(M)$, $u\ne0$. Ceci implique, si $\operatorname{End}(M)$ de dim finie sur $K$, que $\operatorname{End}(M)$ est $\frac12$ simple, et pour tt ss-objet $N\hookrightarrow M$, $\varphi|N$ est une forme de Weil.

\inserted{\uline{\unsure{donc que|dire que} $M$ est semi-simple} s'il est de long. finie (tt ss-objet a un complémentaire)}

3) \deleted{\unsure{On}} Soient $\varphi\in\operatorname{Hom}(M\otimes M,L)$, $\psi\in\operatorname{Hom}(N\otimes N,L)$, $\varphi$ et $\psi$ $\varepsilon$-sym. $(\varepsilon=\pm1)$. Alors $\varphi\oplus\psi$ satisfait à condition de Weil ssi $\varphi$ et $\psi$ y satisfont, et si de plus, pour tt $u:M\unsure{-|\to}N$, on a

% slanted
\margin{\unsure{Généraliser|Généralisation} au cas \deleted{\unsure{où les}} \unsure{d'une cat|d'un corps|un nombre} quelconque de facteurs (\unsure{non cyclotomie|en coefficients|ou sommandes})}

\[
\operatorname{Tr}u^*u>0\qquad\text{si }u\ne0
\]

[$u^*$ calculé pour les formes $\varphi$ et $\psi$ sur $M$ et $N$].

% archive sheet number: 3
\scanpage{5}

On dit alors que les formes $\varphi$ et $\psi$ \deleted{(satisfaisant aux conditions de Weil)} sont \uline{compatibles}.

4) La compatibilité \deleted{des} formes de Weil est une relation \deleted{d'équivalence} réflexive et symétrique. \deleted{\unsure{E|De}} \inserted{\deleted{De façon plus précise, si $\varphi$ est une forme de}} $\varphi,\varphi'$ sur $M$, $\psi,\psi'$ sur $N$ et \deleted{$\chi$ \unsure{sont deux des formes de} Weil, si $\varphi$ est compatible à $\psi$, $\psi$ à $\chi$, et s'il existe une forme} \inserted{si $\varphi,\varphi'$ compatibles, $\psi,\psi'$ compatibles,} alors $\varphi$ et $\psi$ sont compatibles ssi $\varphi'$ et $\psi'$ le sont. Pour le voir, il faut expliciter :

% slanted
\margin{\deleted{\unsure{$\varphi,\psi$, $\varphi$ est} \unsure{compatible} \unsure{ssi $\varphi$ et $\psi$}} relation pas transitive, car si $M,N$ sont tels que $\operatorname{Hom}(M,N)=\operatorname{Hom}(N,M)=0$, $M\ne0$, alors pour deux \deleted{formes $\varphi,\unsure{\psi}$} \inserted{de Weil} $\varphi$ sur $M$, $\psi$ sur $N$, on a $\varphi$ comp. à $\psi$, $\psi$ comp. avec $-\varphi$ mais pas $\varphi$ avec $-\varphi$.}

5) \deleted{Soient} $\varphi$ et $\psi$ deux formes $\varepsilon$-symétriques $M\otimes M\to L$, \deleted{\unsure{on}} $u\mapsto u^*$ de $\operatorname{End}(M)$ associée à $\varphi$, et $\psi$ définie par $\alpha\in\operatorname{End}(M)$, $\psi(x,y)=\varphi(x,\alpha y)$, i.e. $\widetilde\psi=\widetilde\varphi\circ\alpha$. Alors $u^{\operatorname{ad}(\psi)}$ défini par $\psi$ est \deleted{\unsure{donné par}} \inserted{égal à} $\alpha^{-1}u^*\alpha$, \deleted{donc $\varphi$ $\psi$ est compatible avec $\varphi$ ssi $\operatorname{Tr}$} et $u^{\operatorname{ad}(\varphi,\psi)}$ \deleted{\unsure{défini|défini par $\psi$}} égal à $u^*\alpha$. \deleted{Donc}

\[
\begin{aligned}
&\operatorname{Tr}u\,u^{\operatorname{ad}(\varphi,\psi)}>0 &&\text{pour tt }u\ne0\\
&\hspace{2.7em}\Updownarrow\\
&\operatorname{Tr}uu^*\alpha>0 &&\text{pour tt }u\ne0\\
&\hspace{2.7em}\rotatebox{90}{$=$}\\
&\hspace{1.6em}\operatorname{Tr}u^*\alpha u
\end{aligned}
\]

De plus, cette condition implique que \deleted{pour tt} \inserted{$\psi$ est de Weil}, i.e.

\[
\begin{aligned}
&\operatorname{Tr}u\,u^{\operatorname{ad}(\psi,\psi)}>0 &&\text{pour }u\ne0\\
\text{i.e.}\quad &\operatorname{Tr}u\alpha^{-1}u^*\alpha>0 &&\text{pour }u\ne0
\end{aligned}
\]

Donc les formes de Weil sur $M$ compatibles à $\varphi$ correspondent aux \deleted{\unsure{homs}} $\alpha\in\operatorname{End}(M)$ tels que $\operatorname{Tr}u^*u\alpha>0$ pour $u\ne0$. \inserted{\deleted{\unsure{donc}} Donc, pour les formes \inserted{de Weil} sur $M$, la relation de compatibilité est une relation d'équivalence.}

% slanted
\margin{Conséquence : si $\varphi$ et $\psi$ sont de Weil compatibles sur un $=$ $M$, $\varphi\unsure{+\!|\circ|\oplus}\psi$ est de Weil et compatible avec $\varphi$ et $\psi$}

6) Soient \deleted{\unsure{$\varphi_i$}} $(M_i,\varphi_i)$, \deleted{des} avec $\varphi_i:M_i\otimes M_i\to L$, $\varphi_i$ $\varepsilon$-symétriques de Weil ($\varepsilon=\pm1$ fixé). Supposons que pour $i,j$, $\varphi_i$ et $\varphi_j$ compatibles. Alors

% archive sheet number: 4
\scanpage{6}

cette condition est encore remplie si on stabilise en prenant des sommes directes de formes, et des formes induites sur sous-objets.

\inserted{\deleted{\unsure{Plus précis, pour $\varphi$ fixé, on a}} \unsure{Si|de $\sum_I$,} $I$ fini, les formes sur $\coprod M_i$ qui sont compatibles aux $\varphi_i$ sont celles qui sont compatibles \deleted{avec} à $\coprod_i\varphi_i$. Si $\varphi$ de Weil sur $M$, et $N\hookrightarrow M$, les formes sur $N$ comp. à $\varphi$ sont celles compatibles à $\varphi|N$}

7) \deleted{Appelons} Supposons $\mathcal M$ gradué, muni d'un ``objet de Tate'' $K(1)$ de degré $-2$. On appelle \uline{prépolarisation} \deleted{la donn} de $\mathcal M$ la donnée, \deleted{\unsure{d'une}} pour tt $M\in\operatorname{Ob}\mathcal M$ de degré $i$, d'une partie \uline{non vide} \deleted{\unsure{de}} $P(M)\subset\operatorname{Hom}(M\otimes M,K(-i))$, formé de formes $(-1)^i$-sym., \deleted{de Weil,} satisfaisant la condition

\[
\left\{\begin{aligned}
&\varphi\in P(M),\ \psi\in P(N)\quad(M,N\in\operatorname{Ob}\mathcal M^i),\ u:M\to N\quad\text{implique}\\
&\operatorname{Tr}u^{*(\varphi,\psi)}u>0\qquad\text{si }u\ne0.
\end{aligned}\right. \tag*{(P1)}
\]

Cela signifie que pour les $M\in\operatorname{Ob}\mathcal M^i$ ($i$ fixé) les $\varphi\in P(M)$ sont mutuellement compatibles.

\margin{\uline{NB} Si $\exists$ prépolarisation $P$, alors $\mathcal M$ est semi-simple}

Si $P$ et $P'$ sont des prépolarisations, alors $P\subset P'\overset{\text{déf}}{\Longleftrightarrow}\forall M$ de degré $i$, on a $P(M)\subset P'(M)$. Notion de prépolarisations maximales. Grâce à 4) on trouve

\uline{Proposition} Toute prépolarisation $P$ est contenue dans une unique prépolarisation maximale $\widehat P$. On a $\widehat P(M)=$ ensemble des formes de Weil sur $M$ compatibles aux formes de $P(M)$.

\uline{Cor.} a) Si $\varphi\in\widehat P(M)$, $\psi\in\widehat P(N)$, alors $\varphi\oplus\psi\in\widehat P(M\oplus N)$ ($M$ et $N$ de $=$ degré $i$)

b) Si $\varphi\in\widehat P(M)$ \deleted{\unsure{$\varphi\in\widehat P$}} $N\subset M$, alors $\varphi|N\in\widehat P(N)$

c) si $\varphi\in\widehat P(M)$, alors $\check\varphi(\unsure{-i|i})\in\widehat P(\check M\deleted{\unsure{)}}(-i))$

% archive sheet number: 5
\scanpage{7}

8) \deleted{\unsure{Polar|Pour}} La donnée d'une prépolarisation $P$ se décompose en des composantes $P^i$ ($i\in\mathbb Z$). Elles sont à priori tt à fait indépendantes. Une prépolarisation est dite admissible si elle satisfait les conditions

\[
\begin{aligned}
&\varphi\in P(M),\ \psi\in P(N),\ \chi\in P(M\otimes N)\\
&\quad\longrightarrow\unsure{\chi|\varphi}\ \text{compatible avec }\varphi\otimes\psi\quad\text{i.e. dans }\widehat P(M\otimes N)
\end{aligned}\tag*{(P2)}
\]

\[
\begin{aligned}
&\forall i,\quad\text{l'hom can }K(\unsure{1|i})\otimes K(\unsure{1|i})\to K(\unsure{2i|2|1i})\quad\deleted{\text{\unsure{prescrit}}}\\
&\hspace{13em}\text{est dans }\widehat P(K(\unsure{1|i})).
\end{aligned}\tag*{(P3)}
\]

En présence de (P2), (P3) équivaut à

\deleted{$(\mathrm{P}'3)_0$ \inserted{l'hom can} $K(0)\otimes K(0)\to K(0)$ est dans $\widehat P(K(0))$}

\[
\text{Si }\varphi\in P(M),\quad\text{alors }\varphi(i)\in\widehat P(M(i)).\tag*{$(\mathrm P'3)_{\unsure{3|i|1}}$}
\]

\deleted{(Il suffit de le supposer pour $i=1$) \unsure{et}}

Une \uline{polarisation} est une prépolarisation maximale et admissible.

\uline{NB} Si $P$ est une prépolarisation, elle est admissible sss $\widehat P$ est une polarisation. Ceci provient du fait suivant : si $\varphi\in\operatorname{Hom}(M\otimes M,L)$, $\varphi'\in\operatorname{Hom}(M'\otimes M',L')$ ($L,L'$ inv), et si $\varphi,\varphi',\varphi\otimes\varphi'\in\operatorname{Hom}(\underbrace{M\otimes M'}\otimes(M\otimes M'),L\otimes L')$ sont de Weil, alors $\varphi_1$ comp. à $\varphi$ et $\varphi'_1$ comp. à $\varphi'$ implique $\varphi_1\otimes\varphi'_1$ compatible à $\varphi\otimes\varphi'$.

% slanted
\margin{\uline{vérifier}}

9) \deleted{Soit} $P$ une polarisation. Alors on a, pour les formes \inserted{de polarisation}

a) stabilité par $\varphi\oplus\psi$ \inserted{a$'$) \deleted{\unsure{st}} stabilité par $\varphi+\psi$ [résulte de a) et b)]}

b) stabilité par $\varphi|N$

c) stabilité par $\varphi\otimes\psi$

e) stabilité par $\check\varphi$

d) stabilité par $\varphi(i)$

f) \deleted{Contient} $K(0)\otimes K(0)\xrightarrow{\sim}K(0)$ est de \unsure{pol} sss il est défini par $\lambda\in K$, avec $\lambda>0$

\margin{\uline{NB} les conditions b) c) e) f) impliquent déjà que $P$ est une prépolarisation. Les $P$ maximaux satisfaisant b) c) e) f) sont les prépolarisations maximales satisfaisant (P2). Donc les polarisations $\Rightarrow$ systèmes \deleted{\unsure{maximaux}} satisfaisant b) \unsure{c)} d) e) f) et \uline{maximaux} (\uline{NB} a) en résulte)}

% archive sheet number: 6
\scanpage{8}

10) Soit $M_\alpha$ un système d' \deleted{$\otimes$-générateurs de} \inserted{\unsure{élém$^{\rm ts}$|élém.} \deleted{\unsure{de}} hom. de} $\mathcal M$, tel que avec $K(1)$ et $K(-1)$ ils engendrent $\mathcal M$ comme $\otimes$-catégorie. Alors une polarisation est connue quand, pour tt $\alpha$, on connaît \uline{une} forme de polarisation \inserted{$\varphi_\alpha$} de $M_\alpha$.

On peut supposer d'ailleurs que les $M_\alpha$ sont tous de degré $0$ ou $1$, et si $\mathcal M$ est à engendrement fini, que le syst. est réduit à un $M^0$ et un $M^1$. \deleted{Ils correspondent à des \unsure{=}}

11) \uline{Remarque 1} \deleted{Les polarisati} Pour $M,L$ fixés, les classes d'équivalence de formes de Weil \deleted{\unsure{sur}} $M\otimes M\to L$, \deleted{pour} pour $\mathcal M$ ou $\mathcal M\otimes_K\mathbb R$, sont les $=$. La compatibilité des formes de Weil est la $=$, au sens \deleted{$\mathcal M\otimes_K\mathbb R$} $\mathcal M$ ou $\mathcal M\otimes_K\mathbb R$. Ils correspondent à certaines composantes connexes de l'ouvert de $\operatorname{Hom}(M_{\mathbb R}\otimes M_{\mathbb R},L_{\mathbb R})$ formé des formes $\varepsilon$-sym. \uline{non dégénérées}, et \inserted{sont \uline{donc} \uline{en nb fini}.}

\uline{Remarque 2} Les prépolarisations maximales de $\mathcal M^i$ ou $\mathcal M_{\mathbb R}^i$ sont en corr. biunivoque. De $=$ donc pour $\mathcal M$ et $\mathcal M_{\mathbb R}$. Si $P,P_{\mathbb R}$ se correspondent, alors $P$ est une polarisation sss $P_{\mathbb R}$ en est une. Si $\mathcal M$ est à engendrement fini, \deleted{\unsure{on trouve|c) Prouver}} l'ens des polarisations possibles de $\mathcal M$ est fini.

On est ainsi ramené à étudier les polarisations des $\otimes$-catégories \uline{sur} $\mathbb R$. De plus

\uline{Remarque 3} Soit $\mathcal M=\varinjlim\mathcal M_{\unsure{(\alpha)|\alpha}}$\deleted{\unsure{cat}} \inserted{sous-$\otimes$-catégories famille filtrante croissante}, alors l'ens $P(\mathcal M)$ des polarisations de $\mathcal M$ est donné par $P(\mathcal M)=\varprojlim_\alpha\mathcal M(\alpha)$ \deleted{\unsure{donc} $P(\mathcal M)$} Donc si les $\mathcal M(\alpha)$ \deleted{sont} de t.f., on

% archive sheet number: 7
\scanpage{9}

voit que $P(\mathcal M)\ne\varnothing\Longleftrightarrow\forall\alpha,\ P(M_\alpha)\ne\varnothing$. On est ainsi ramené au cas de \uline{$\otimes$-catégories semi-simples de type fini} \uline{sur} $\mathbb R$ : étudier l'ens. des polarisations d'une telle $\mathcal M$.

12) Soit $\mathcal M$ avec polarisation $P$ (sur $\mathbb R$ si on veut). Soit $M\in\operatorname{Ob}\mathcal M_{\mathbb C}^i$, i.e. un objet \inserted{$M_0$} de $\mathcal M^i$, avec action de $\mathbb C$ dessus\unsure{\,\textit{lui}|\,en} (i.e. endomorphisme $J_M$ de carré $-1$). Soit $\overline M$ avec l'action conjuguée ($J_{\overline M}=-J_M=J_M^{-1}$). Un hom

\[
M\otimes_{\mathbb C}\overline M\xrightarrow{\varphi}\mathbb C(-i)=\mathbb R(-i)\otimes_{\mathbb R}\mathbb C
\]

équivaut à un hom

\[
(M\otimes_{\mathbb C}\overline M)_{\underset{\substack{\inserted{\text{motif réel}}\\\text{sous-jacent}}}{0}}\longrightarrow\mathbb R(-i)
\]

i.e un hom

\[
M_0\otimes M_0\xrightarrow{\varphi_0}\mathbb R(-i)
\]

\[
\begin{aligned}
\text{satisfaisant}\quad\varphi_0(J_Mx,y)&=\varphi_0(\deleted{\unsure{X}}x,J_{\overline M}y)\\
&=-\varphi_0(x,J_My)\qquad\text{i.e.}\\
\varphi_0(J_Mx,J_My)&=\varphi_0(x,y)
\end{aligned}
\]

Pour que $\varphi$ soit \deleted{\unsure{soit}} ``hermitienne'' resp. ``antihermitienne'' il faut \unsure{s.|et s.|si} que $\varphi_0$ soit sym. resp. alternée [la symétrie hermitienne en $\varphi$ correspond à la symétrie ordinaire en $\varphi_0$]. On dit que $\varphi$ est une \uline{polarisation} \inserted{(pour $P$)} si $\varphi_0$ \deleted{est} est une (pour $P$). Cela implique $\varphi$ non dégénérée, i.e. $\widetilde\varphi:\check M\simeq\overline M(-i)$

% archive sheet number: 8
\scanpage{10}

\uline{Il y a tjrs une forme de polarisation $P$ sur} $M$, i.e. \deleted{en effet, il suffit de} une polarisation $\varphi_0$ de $M_0$ invariante par $J_M$. En effet, si $\psi_0$ est une polarisation de $M_0$, on pose

\[
\begin{aligned}
\varphi_0(x,y)&=\psi_0(x,y)+\psi_0(J_Mx,J_My)\\
\text{donc}\quad\varphi_0(J_Mx,J_My)&=\psi_0(J_Mx,J_My)+\psi_0(-x,-y)=\varphi_0(x,y).
\end{aligned}
\]

Par suite, on a $\check M\simeq\overline M(-i)$ pour tt $M\in\operatorname{ob}\mathcal M_{\mathbb C}^i$ ).

13). On voit que pour une $\mathcal M$ donnée \inserted{sur $\mathbb R$}, on a des implications

\diagram{\begin{tikzcd}[ampersand replacement=\&,row sep=small,column sep=small]
{}\arrow[d,Rightarrow] \& \alpha)\quad\mathcal M\text{ admet une polarisation} \& {}\\
{}\arrow[d,Rightarrow] \& \beta)\quad\forall M\in\operatorname{Ob}\mathcal M_{\mathbb C}^{i},\ \text{on a }\check M\simeq\overline M(-i), \& \inserted{\mathcal M\ \frac12\text{ simple}}\\
{} \& \gamma)\quad\forall M\in\operatorname{Ob}\mathcal M_{\deleted{\unsure{\mathbb R|\beta}}}^{i},\ \text{on a }\check M\simeq M(-i), \& \inserted{\mathcal M\ \frac12\text{ simple}}
\end{tikzcd}}

($\beta\Rightarrow\gamma$ provient du fait que $M_{\mathbb C}\simeq N_{\mathbb C}$ implique $M\simeq N$).

\margin{\uline{Lemme}}

Les implications précédentes \deleted{sont} \unsure{s'illustrent|sont strictes} chez \inserted{au moins}, $\mathcal M=\operatorname{Rep}(G)$, $G$ groupe algébrique \inserted{affine \deleted{\unsure{connexe} \illegible{5}}} \deleted{\unsure{simple}} sur $\mathbb R$. On peut avoir $\beta$) \deleted{\unsure{et non|sans avoir}} sans avoir $\alpha$)

\margin{?}

[$\deleted{\operatorname{SO}(\varphi)}$ \deleted{\unsure{de}} $\frac{\operatorname{SO}(\varphi)}{\text{Centre}}$, $\varphi$ forme à un nb pair de variables, avec un nombre impair de carrés négatifs], et $\gamma$) sans avoir $\beta$) [$\ldotp\ldotp\ldotp\ldotp$\quad ?\quad$\ldotp\ldotp\ldotp\ldotp\ldotp$]

% archive sheet number: 9
\scanpage{11}

14) Supposons $\mathcal M$ semi-simple \inserted{sur $K\subset\mathbb R$} \unsure{engendrée|engendré} par motifs $M_1\ldots M_r$ de degrés $d_1\ldots d_r$, \inserted{et $K(1)$ \unsure{resp.|,} $K(-1)$} \deleted{\unsure{de}} \deleted{\unsure{objets}}. Pour tt $i$, soit $\varphi_i$ forme $(-1)^{d_i}$-sym. sur $M_i$, qui soit \deleted{\unsure{déf.} $\unsure{M_i}$} forme de Weil. \deleted{\unsure{et}} Pour qu'il existe une polarisation de $\mathcal M$ \deleted{définie} pour laquelle les $\varphi_i$ sont des formes de polarisation, il \deleted{\unsure{s}} faut qu'on ait la condition suivante :

a) Pour deux systèmes d'entiers $\underline m=(m_1,\ldots,m_{r+1}),\underline m'=(m'_1,\ldots,m'_{r+1})$ $\geq0$, avec $\overbrace{m_{r+1}+\sum_1^r m_i d_i}^{\delta(\underline m)}=\overbrace{m'_{r+1}+\sum m'_i d'_i}^{\delta(\underline m')}$, posant

\[
M_{\underline m}=\left(\bigotimes_{1\leq i\leq r}\left(\bigotimes^{m_i} M_i\right)\right)(\underset{\unsure{2|?}}{\unsure{+|-}}m_{r+1}),\quad
\varphi_{\underline m}=\left(\bigotimes_{1\leq i\leq r}\left(\underset{\deleted{\unsure{\varphi}}}{\bigotimes^{m_i}}\varphi_i\right)\right)(\underset{\unsure{2|?}}{\unsure{-|+}}m_{r+1}),\ \deleted{\unsure{\alpha}}
\]

on a, pour tout $u\in\operatorname{Hom}(M_{\underline m},M_{\underline m'})$, la relation

\[
\operatorname{Tr}u^{*(\varphi_{\underline m},\varphi_{\underline m'})}u\geq0
\]

Je dis que cette condition est aussi suffisante. \deleted{\unsure{Si elle est}} \inserted{Supposons celà en effet} \unsure{remplie|rempli}, soit $M\in\operatorname{Ob}\mathcal M^k$, \inserted{$M\ne0$.} par hyp. $\exists\ \underline m$ avec un $\alpha:M\hookrightarrow M_{\underline m}$ ($k=\deleted{\text{deg}}\sum m_i d_i$), soit $P(M)$ l'ens des formes \deleted{qui} sur $M$ qui sont compatibles avec $\varphi_{\underline m}|M$. Prouvons que c'est indépendant de $\underline m,\alpha$.

En effet, si on a aussi $\alpha':M\unsure{\subset|\hookrightarrow}M_{\underline m'}$, comme $\varphi_{\underline m}$ et $\varphi_{\underline m'}$ sont compatibles par hyp, il en est de $=$ de $\varphi_{\underline m}|M$ et $\varphi_{\underline m'}|M$. On voit que ces $P(M)$ forment une prépolarisation maximale, \deleted{\unsure{et}} à laquelle appartiennent \inserted{évidemment} les $\varphi_i$. Si $\varphi\in P(M),\varphi'\in P(M')$, \deleted{\unsure{qui}} on a $\varphi\sim\varphi_{\underline m}|M$, $\varphi'\sim\varphi_{\underline m'}|M'$, donc $\varphi\otimes\varphi'\sim\varphi_{\underline m}\otimes\varphi_{\underline m'}|M\otimes\unsure{M'|M)}=\varphi_{(\underline m+\underline m')}|M\otimes M'$, \deleted{\unsure{donc}}

% archive sheet number: 10
\scanpage{12}

donc $\varphi\otimes\varphi'\in P(M\otimes M')$. Enfin, supposons $\varphi\in P(M)$, i.e. $\varphi\sim\varphi_{\underline m}|M$, alors $\varphi(i)\sim\varphi_{\underline m}(i)|M(i)\unsure{;|,}=\varphi_{(m_1,\ldots,m_r,m_{r+1}-i)}|M(i)$, donc $\varphi(i)\in P(M(i))$.

15) Soit $G$ groupe algébrique sur $\mathbb R$, \deleted{$C\in$} avec

\[
G_m\xrightarrow{i}G\xrightarrow{\varepsilon}G_m,\quad\unsure{y|\gamma|\eta}=i(-1),\quad\text{et soit }C\in G(\mathbb R)
\]

avec $C^2=y$. On prend $\mathcal M=\operatorname{Rep}(G)$

\deleted{1) Soit $M$} a) Si $M\in\operatorname{Ob}\mathcal M^i$ i.e. \inserted{muni d'une} repr de \deleted{degré} poids $i$, \deleted{\unsure{Obs.}} si $\varphi:M\otimes M\unsure{-|\to}\mathbb R(-i)$ est telle que $\varphi(x,Cy)$ est symétrique définie positive, alors $\varphi$ est une forme de Weil, et si on a de $=$ $\psi$ \deleted{\unsure{$\longmapsto$}} $:N\otimes N-\mathbb R(-i)$, $\psi(x,Cy)$ déf. positive, alors $\varphi$ et $\psi$ sont \uline{compatibles}.

\margin{$\varphi$ est dite une forme de $C$-polarisation \deleted{$\varphi$ est}}

% slanted
\margin{vérifions d'abord que $\varphi$ est \deleted{\unsure{sym}} $(-1)^i$-sym.}

En effet, soit $u:M\longrightarrow N$ compatible avec l'action de $G$ : Considérons $u^{*(\varphi,\psi)}$, et comparons avec $u^{*(\varphi^C,\psi^C)}$ ($\varphi^C(x,y)=\varphi(x,Cy)$). On aura

\[
\begin{aligned}
u^{*(\varphi^C,\psi^C)}&=C_M^{-1}u^{*(\varphi,\psi)}C_N\qquad\text{donc}\\
u^{*(\varphi,\psi)}&=C_M\underbrace{u^{*(\varphi^C,\psi^C)}}_{u'}C_N^{-1},\qquad\text{donc}\\
\operatorname{Tr}uu^{*(\varphi,\psi)}&=\operatorname{Tr}uC_Mu'C_N^{-1}
=\operatorname{Tr}\underbrace{C_N^{-1}u}_{uC_M^{-1}}C_Mu'=\operatorname{Tr}uu'\geq0
\end{aligned}
\]

% slanted
\margin{\unsure{dis q. tjrs la même|dis q. tjrs le même} (sans condition de positivité) $\operatorname{Tr}uu^{*(\varphi^{\unsure{*|C}},\psi)}\deleted{\unsure{0}}=\operatorname{Tr}uu^{*(\varphi^C,\psi^C)}$}

\deleted{Or $C_M$ et $C_N$ sont orthogonaux pour les formes déf positives $\varphi^C,\psi^C$ (en fait, les $g\in G(\mathbb R)$ qui dans $\operatorname{Centr}(C)$ opèrent de façon orthogonale), on a donc $C_N^{-1}=C$}

b) Si $\varphi,\psi$ sont des $C$-polarisations, $\varphi\otimes\psi$ aussi, si $\varphi$ une $C$-polarisation, alors $\varphi(i)$ aussi, \inserted{\uline{pourvu que} $\varepsilon(C)=1$ (et non $-1$ !)} la forme \unsure{standart|standard} $\mathbb R\otimes\mathbb R-\mathbb R$ est une $C$-polarisation.

% archive sheet number: 11
\scanpage{13}

Pour \deleted{que} une $\deleted{\unsure{\varphi|C}}C$-polarisation donnée, les formes $\varphi'$ sur le $=$ $M$ qui sont des $C$-polarisations, sont \uline{exactement} celles qui sont compatibles à $\varphi$.

c) Donc : Supposons $\varepsilon(C)=1$. Soit, pour tt $M$ donné, $P(M)$ l'ens de ses $C$-polarisations. Alors $M\mapsto P(M)$ est une polarisation de $\mathcal M$, \deleted{\unsure{à|et}} cela près que les $P(M)$ peuvent ne pas être tous $\ne\varnothing$.

16) \uline{Questions} Toute polarisation de \deleted{\unsure{$\mathcal M$}} $\mathcal M=\operatorname{Rep}(G)$ peut-elle se décrire par un $C$ comme dessus ($C^2=\unsure{y|\gamma|\eta},\ \varepsilon(C)=1$) ? Est-ce vrai si $G$ connexe ? Cas $G'$ \inserted{$=\operatorname{Ker}\varepsilon$} compact ?

\margin{$?$}

\uline{NB} Un tel $C$ \inserted{mod centre $Z$} est uniquement déterminé mod \deleted{la} \deleted{\unsure{centralité}} mod conjugaison, \deleted{\unsure{pour éléments}} du moins si $G$ est connexe. Donc \deleted{\unsure{l'on|bien}} l'indétermination de $C$ \inserted{dans ce cas} est modulo points d'ordre $2$ de $Z(\mathbb R)$ et conjugaison.

\uline{Re NB} Pour qu'il existe un $C$ comme dessus fournissant une polarisation, il \inserted{cf. 14} suffit, en termes du groupe $G'=\operatorname{Ker}\varepsilon$ muni de $y=i(-1)$, qu'il existe un $C\in G'(\mathbb R)$ tel que \uline{\unsure{toute|toutes}} \unsure{repr|reprs|rep.} \inserted{de $G'$} sur laquelle $y$ induit $\pm1$ admette une $C$-polarisation, i.e. une forme \inserted{inv.} $\varphi$ telle que $\varphi(x,Cy)$ soit sym. déf $>0$.

On dit que le groupe \inserted{alg affine} réel $G'$ est \uline{hodgien} s'il existe un tel $C$ \deleted{d'ordre 4} \inserted{avec $C^4=1$} ($y=C^2$ n'étant pas fixé à priori).

% archive sheet number: 12
\scanpage{14}

17) Soit $C\in G'(\mathbb R)$ \deleted{\unsure{$\varepsilon$}}, avec $C^2=\unsure{y|\gamma|\eta}$ $(=i(-1))$, définissant une polarisation de $\mathcal M=\operatorname{Rep}(G)$, -- d'un $G$ réductif !

Alors $\overbrace{\operatorname{Centr}_{G'}(C)}^{K}(\mathbb R)$ est un \uline{compact} maximal de $G'(\mathbb R)$. En effet, on \deleted{\unsure{voit tout de suite que}} \deleted{\unsure{$\varphi$}} aura \deleted{$G'\subset\operatorname{Aut}_{\mathbb R}(M^0)\times\operatorname{Aut}_{\mathbb R}(M^1)$, et} $G'\hookrightarrow\prod_\alpha\operatorname{Aut}_{\mathbb R}(M_\alpha)$, par des repr. convenables (au nb de 2 si on veut) de $G'$, \deleted{de} avec $y$ y opérant par $\pm1$. Si $\varphi_\alpha$ est une forme de $C$-polarisation sur $M_\alpha$, alors les $g\in\operatorname{Centr}_{G'}(C)$ laissent invariant $\varphi_\alpha^C$, donc

\[
\deleted{\unsure{\operatorname{Centr}|\operatorname{Cent}}}\quad K\subset\prod_\alpha\underbrace{\operatorname{Aut}_{\mathbb R}(M_\alpha,\varphi_\alpha^C)}_{\substack{\text{groupes réels}\\\text{compacts}}}\ ,
\]

donc $K$ est \inserted{en pts réels} compact. Prouvons que \deleted{il est} $K(\mathbb R)$ est compact \uline{maximal}. Soit \deleted{donc} $\underline{K'}\subset G'(\mathbb R)$, $K'\supset K(\mathbb R)$. On sait que $\underline{K'}$ est algébrique, $\underline{K'}=K'(\mathbb R)$, $K'$ ss-groupe algébrique de $G'$ contenant $C$. Quitte à remplacer $G'$ par $K'$, on peut supposer $G'$ compact, et il faut prouver qu'alors $C$ est central.

Passons le d'abord dans $G'/\{1,y\}$, i.e. supposons d'abord $y=1$. Soit une repr. fidèle de $G'$ dans un $V$, \deleted{\unsure{avec forme de pol. sur} $V$} \deleted{Alors $C$ \unsure{définit}} \inserted{avec forme de $C$-pol. $\varphi$ sur $V$.} $V=V^++V^-$

\deleted{($C\mid V^+=\mathrm{id}$, $C\mid V^-=-\mathrm{id}$) et} avec $(\varphi\mid V^+)$ \deleted{est} $>0$, $\underset{\substack{\Vert\\-\varphi^-}}{(\varphi\mid V^-)}<0$.

\deleted{\unsure{$\varepsilon$}} $G'$ étant un sous-groupe \unsure{Cpt|cpt} de $O(\varphi)$, on peut par conjugaison \inserted{par $u\in O(\varphi)(\mathbb R)$} l'envoyer dans le compact maximal $O(V^+)\times O(V^-)$ \inserted{de $O(\varphi)$}, donc on peut supposer que $G\subset O(V^+)\times O(V^-)$. Donc $C\in G(\mathbb R)$ est de la forme $C^+\oplus C^-$, \deleted{\unsure{et alors}} Comme $\varphi^+(x,C^+y)$ et

% archive sheet number: 13
\scanpage{15}

$\varphi^-(x,C^-y)$ sur $V^+$ resp. $V^-$ sont définies positives, $C^+$ et $C^-$ étant des \deleted{\unsure{invol}} involutions orthogonales, on conclut que $C^+=\mathrm{id}$, $C^-=-\mathrm{id}$, donc $C$ est dans le centre de \deleted{$O(\varphi)$, \unsure{$\varphi$}} $O(V^+)\times O(V^-)$ donc dans le centre de \unsure{$G$|$G'$}.

Cas $\unsure{y|\gamma|\eta}\ne1$. Alors $gCg^{-1}\in\{C,yC\}$, donc $G'$ \deleted{\unsure{général} ($y$ pas néc. $=1$). On voit déjà que} \deleted{\unsure{est un morphisme} $G\to$} \deleted{$g\mapsto[C,g]=CgC^{-1}g^{-1}\in\{1,y\}$ pour tt $g\in G'(\mathbb R)$.}

\deleted{Soit une représentation \inserted{simple non nulle} de $G$ telle que $y$ y opère par $-\mathrm{id}$. il suffit de prouver que dans l'image, $C$ devient central. Donc on peut supposer $G\subset\operatorname{Aut}_{\mathbb R}(M)$, $y$ opérant par $-\mathrm{id}$. Donc $C$ définit sur $M$ structure complexe, et $G\subset\operatorname{Aut}_{\mathbb C}(M)$, La forme $\varphi$ de $C$-polarisation satisfait à}

opère sur $\{C,yC\}$, \deleted{\unsure{et il y a un sous-groupe}} soit $G''$ le stabilisateur de \unsure{$G$|$C$}, alors $G''$ est d'indice 1 ou 2, \deleted{\unsure{On obtient}} \inserted{et} \deleted{\unsure{on voit que l'on a}} $gCg^{-1}=\varepsilon(g)C$\quad $(\varepsilon:G\longrightarrow\mu_2)$

\[
\varphi^C(gx,gy)=\varphi(gx,Cgy)=\varepsilon(g)\,\varphi(gx,gCy)=\varepsilon(g)\varphi^C(x,y).
\]

\deleted{Donc $G'\subset\operatorname{Sp}^{\unsure{1|\prime}}(\varphi^C)$}

Si on avait $\varepsilon(g)=-1$, on aurait $\varphi^C(gx,gx)=-\varphi^C(x,x)$, et si $x\ne0$ les deux mbs seraient resp $>0$ et $<0$, absurde.

% archive sheet number: 14
\scanpage{16}

18) Supposons que $C,C'$ \deleted{\unsure{$\in$}} \inserted{$(C^2=C'^2=\unsure{y|\gamma|\eta})$} définissent des polarisations de \deleted{\unsure{$G$}} \inserted{$\mathcal M=\operatorname{Rep}(G)$}.

Quand celles-ci sont-elles \uline{égales} ? Il \uline{suffit} qu'on ait $C'=gCg^{-1}$, avec $g\in G(\mathbb R)$, car alors pour $\varphi$ forme invariante donnée, on aura

\[
\varphi^{C'}(x,y)=\varphi(x,gCg^{-1}y)=\varphi(g^{-1}x,Cg^{-1}y)=\varphi^C(g^{-1}x,g^{-1}y);
\]

Cette condition est-elle aussi \uline{nécessaire} ? \unsure{Oui|On}, \unsure{si on suppose} \inserted{$G'$ \unsure{connexe|connexes}} et on trouve même \unsure{$=$} la conjugaison par un élément de $G'(\mathbb R)^0$ !

% slanted
\margin{\deleted{\unsure{Quand}} \uline{NB} $G'$ connexe $\unsure{\Longleftrightarrow|\Longleftarrow}$ $\unsure{G|C}$ connexe (grâce à l'existence du $C$ \unsure{avec|ou} $C^2=y$)}

Tout d'abord, la conjugaison des gpes maximaux dans un gpe de Lie connexe montre que, \deleted{\unsure{les}} si $C,C'$ sont deux élémts de $G'$ qui sont hodgiens (sans avoir même $C^2=C'^2$ !) il existe $g\in G'(\mathbb R)^0$ tel que

\[
gCg^{-1}C'^{-1}\in Z'(\deleted{\unsure{\mathbb R}}\mathbb R)\qquad (\unsure{Z|Z'}=\text{Centre de }G')
\]

[ \uline{NB} $(G'/Z')(\mathbb R)^0\longleftarrow G'(\mathbb R)^0$ est surjectif ! ].

Donc on aura $C'=zgCg^{-1}$, $g\in G'(\mathbb R)^0$, $z\in Z'(\mathbb R)$. Si $C'^2=C^2$, on aura d'ailleurs $z^2=1$ i.e. $z\in {}_2Z'(\mathbb R)$.

\deleted{On \unsure{se}} On peut donc supposer $C'=z\deleted{\unsure{g}}C$, $z\in {}_2Z(\mathbb R)$, et on est ramené à prouver que si $z\ne1$, alors $C'=zC$ et $C$ donnent des polarisations diff.

Or si $z\ne1$, il y aura une repr. de $G$ dans un $V\ne0$, où $z$ opère par $-\mathrm{id}_V$. Si $\varphi$ est une forme $C$-polarisante, $\varphi^C>0$, or $\varphi^{C'}=\varphi^{zC}=-\varphi^C<0$, absurde !

Donc on trouve que l'ensemble des polarisations sur $\mathcal M=\operatorname{Rep}(G)$ qui sont définissables par un $C$ (i.e. qui sont ``hodgiennes'') forme un \uline{torseur} sous $ {}_2Z(\mathbb R)$ : Si \deleted{\unsure{$\varphi$ est une}} $P$ est une polarisation sur $\mathcal M$ et $z\in {}_2Z(\mathbb R)$, alors $z.P$ est la polarisation

% archive sheet number: 15
